% SOURCE: Chinese PDF viewer pp.127--132 (upper part).
\section{Lippmann--Schwinger Equation}
\subsection*{1) Lippmann--Schwinger equation in Hilbert space}
Let
\begin{equation}
\op H=\op H_0+\op V,\qquad \op H_0=\frac{\op{\vect p}^{\,2}}{2m}.
\end{equation}
The incident plane wave is $\ket\varphi$, with
\begin{equation}
\op H_0\ket\varphi=E\ket\varphi,\qquad E=\frac{p^2}{2m}=\frac{\hbar^2 k^2}{2m}.
\end{equation}
The scattered wave is $\ket\psi$. In elastic scattering the energy is unchanged:
\begin{equation}
\op H\ket\psi=(\op H_0+\op V)\ket\psi=E\ket\psi,\qquad
(E-\op H_0)\ket\psi=\op V\ket\psi.
\end{equation}
The formal solution satisfying the initial condition $\lim_{V\to0}\ket\psi=\ket\varphi$ is
\begin{equation}
\ket\psi=\ket\varphi+\frac{1}{E-\op H_0}\op V\ket\psi.
\end{equation}
However, the operator $(E-\op H_0)^{-1}$ diverges when acting on eigenstates of $\op H_0$. To avoid the divergence, extend the energy into the complex plane and write the formal solution as
\begin{equation}
\boxed{\ket{\psi^\pm}=\ket\varphi+
\frac{1}{E-\op H_0\pm\ii\epsilon}\op V\ket{\psi^\pm},\quad\epsilon\to0^+.}
\end{equation}
This is the Lippmann--Schwinger equation in Hilbert space.

Enter the coordinate representation:
\begin{equation}
\braket{\vect r}{\psi^\pm}=\braket{\vect r}{\varphi}
+\int\dd^3r'\,\mel{\vect r}{\frac1{E-\op H_0\pm\ii\epsilon}}{\vect r'}
\mel{\vect r'}{\op V}{\psi^\pm}.
\end{equation}
Because
\begin{equation}\begin{aligned}
\mel{\vect r}{\frac{1}{E-\op H_0\pm\ii\epsilon}}{\vect r'}
&=\int\dd^3p'\dd^3p''\,\braket{\vect r}{\vect p'}
\mel{\vect p'}{\frac1{E-\op H_0\pm\ii\epsilon}}{\vect p''}
\braket{\vect p''}{\vect r'},\\
\mel{\vect p'}{\frac1{E-\op H_0\pm\ii\epsilon}}{\vect p''}
&=\frac{\delta^{(3)}(\vect p'-\vect p'')}{E-E'\pm\ii\epsilon},\\
\braket{\vect r}{\vect p'}&=\frac{e^{\ii\vect p'\cdot\vect r/\hbar}}{(2\pi\hbar)^{3/2}},\qquad
\braket{\vect p''}{\vect r'}=\frac{e^{-\ii\vect p''\cdot\vect r'/\hbar}}{(2\pi\hbar)^{3/2}},
\end{aligned}\end{equation}
we have
\begin{equation}\begin{aligned}
\mel{\vect r}{\frac1{E-\op H_0\pm\ii\epsilon}}{\vect r'}
&=\int\frac{\dd^3p'}{(2\pi\hbar)^3}
\frac{e^{\ii\vect p'\cdot(\vect r-\vect r')/\hbar}}{E-E'\pm\ii\epsilon}\\
&=\int\frac{\dd^3k'}{(2\pi)^3}
\frac{e^{\ii\vect k'\cdot(\vect r-\vect r')}}{E-E'\pm\ii\epsilon}.
\end{aligned}\end{equation}
By the same contour deformation used to define the Green function, moving a pole slightly upward or downward makes the integration path straight. With $E=\hbar^2k^2/(2m)$,
\begin{equation}\begin{aligned}
G^\pm(\vect r-\vect r')
&=\int\frac{\dd^3k'}{(2\pi)^3}
\frac{e^{\ii\vect k'\cdot(\vect r-\vect r')}}{k^2-k'^2\pm\ii0},\\
\mel{\vect r}{\frac1{E-\op H_0\pm\ii\epsilon}}{\vect r'}
&=\frac{2m}{\hbar^2}G^\pm(\vect r-\vect r').
\end{aligned}\end{equation}
Hence the Lippmann--Schwinger equation in the coordinate representation becomes
\begin{equation}
\braket{\vect r}{\psi^\pm}=\braket{\vect r}{\varphi}
+\frac{2m}{\hbar^2}\int\dd^3r'\,G^\pm(\vect r-\vect r')
\mel{\vect r'}{\op V}{\psi^\pm}.
\end{equation}
Since
\begin{equation}\begin{aligned}
\mel{\vect r'}{\op V}{\psi^\pm}
&=\int\dd^3r''\,\mel{\vect r'}{\op V}{\vect r''}
\braket{\vect r''}{\psi^\pm}\\
&=\int\dd^3r''\,V(\vect r')\delta^{(3)}(\vect r'-\vect r'')
\braket{\vect r''}{\psi^\pm}
=V(\vect r')\braket{\vect r'}{\psi^\pm},
\end{aligned}\end{equation}
we obtain
\begin{equation}
\boxed{\braket{\vect r}{\psi^\pm}=
\braket{\vect r}{\varphi}+\int\dd^3r'\,
U(\vect r')G^\pm(\vect r-\vect r')\braket{\vect r'}{\psi^\pm}.}
\end{equation}
This is precisely the integral equation corresponding to the coordinate-space Schr\"odinger equation obtained in the preceding section by the Green-function method. It is the coordinate representation of the Lippmann--Schwinger equation.

\subsection*{2) Dyson equation}
In the general Lippmann--Schwinger equation,
\begin{equation}
\ket{\psi^\pm}=\ket\varphi+\frac1{E-\op H_0\pm\ii\epsilon}\op V\ket{\psi^\pm},
\end{equation}
let
\begin{equation}
\op G_0^\pm=\frac1{E-\op H_0\pm\ii\epsilon}.
\end{equation}
The coordinate representation $\mel{\vect r}{\op G_0^\pm}{\vect r'}$ is a Green function, and thus $\op G_0^\pm$ is called the \emph{free Green operator}, related to $\op H_0$.

The Lippmann--Schwinger equation can also be rewritten as
\begin{equation}\begin{aligned}
(E-\op H_0\pm\ii\epsilon)\ket{\psi^\pm}
&=(E-\op H_0\pm\ii\epsilon)\ket\varphi+\op V\ket{\psi^\pm},\\
(E-\op H\pm\ii\epsilon)\ket{\psi^\pm}
&=(E-\op H\pm\ii\epsilon)\ket\varphi+\op V\ket\varphi,\\
\ket{\psi^\pm}&=\ket\varphi+
\frac1{E-\op H\pm\ii\epsilon}\op V\ket\varphi.
\end{aligned}\end{equation}
Note that $\ket\varphi$ is an eigenstate of $\op H_0$, not of $\op H$. Define
\begin{equation}
\op G^\pm=\frac1{E-\op H\pm\ii\epsilon}
\end{equation}
as the \emph{full Green operator}, related to $\op H$.

From
\begin{equation}\begin{aligned}
(E-\op H\pm\ii\epsilon)\op G^\pm&=\mathbf1,\\
(E-\op H_0\pm\ii\epsilon)\op G^\pm-\op V\op G^\pm&=\mathbf1,
\end{aligned}\end{equation}
we have
\begin{equation}
\boxed{\op G^\pm=\op G_0^\pm+\op G_0^\pm\op V\op G^\pm}.
\end{equation}
This is the \emph{Dyson equation}, which describes the relation between the free and full Green operators. It has important applications in quantum statistical mechanics and quantum field theory.

\subsection*{3) Scattering amplitude}
Consider the state $\ket{\psi^+}$ associated with the retarded Green function $G^+$. The coordinate-space Lippmann--Schwinger equation gives
\begin{equation}
\psi^+(\vect r)=A e^{\ii\vect k\cdot\vect r}
-\frac1{4\pi}\int\dd^3r'\,U(\vect r')\psi^+(\vect r')
\frac{e^{\ii k|\vect r-\vect r'|}}{|\vect r-\vect r'|}.
\end{equation}
Because the effective region of the potential $V(\vect r')$ is finite, when $r\to\infty$,
\begin{equation}\begin{aligned}
|\vect r'|&\ll r,\\
|\vect r-\vect r'|&=\sqrt{(\vect r-\vect r')^2}
=\sqrt{r^2+r'^2-2\vect r\cdot\vect r'}\\
&\simeq r\left(1-2\frac{\vect r\cdot\vect r'}{r^2}\right)^{1/2}
\simeq r-\frac{\vect r\cdot\vect r'}r.
\end{aligned}\end{equation}
Consequently,
\begin{equation}
\psi^+(\vect r)\xrightarrow{r\to\infty}
Ae^{\ii\vect k\cdot\vect r}
-\frac1{4\pi}\frac{2m}{\hbar^2}\frac{e^{\ii kr}}r
\int\dd^3r'\,e^{-\ii\vect k'\cdot\vect r'}V(\vect r')\psi^+(\vect r').
\end{equation}
In the phase, $|\vect r-\vect r'|$ is retained to first order; in the denominator only the zeroth-order term is kept. The scattered wavevector is $\vect k'=k\vect r/r$, with $|\vect k'|=k$.

\begin{center}
\begin{tikzpicture}[>=Latex,scale=.9]
 \coordinate (O) at (0,0);\coordinate (K) at (3.1,0);\coordinate (Kp) at (2.45,2.15);
 \draw[->,thick](O)--(K)node[midway,below]{$\vect k$};
 \draw[->,thick](O)--(Kp)node[midway,above left]{$\vect k'$};
 \draw[->,thick](K)--(Kp)node[midway,right]{$\vect q=\vect k'-\vect k$};
 \draw (0.65,0)arc(0:41:0.65);\node at (0.82,.28){$\theta$};
\end{tikzpicture}
\end{center}
Comparing with the standard asymptotic solution
\begin{equation}
\psi(\vect r)=Ae^{\ii\vect k\cdot\vect r}+Af(\theta,\varphi)\frac{e^{\ii kr}}r,
\end{equation}
we obtain the scattering amplitude
\begin{equation}\begin{aligned}
f(\theta,\varphi)
&=-\frac{(2\pi)^3}{4\pi}\frac{2m}{\hbar^2}
\int\dd^3r'\,\frac{e^{-\ii\vect k'\cdot\vect r'}}{(2\pi)^{3/2}}
V(\vect r')\psi^+(\vect r')\\
&=-\frac{(2\pi)^3}{4\pi}\frac{2m}{\hbar^2}
\int\dd^3r'\,\braket{\vect k'}{\vect r'}V(\vect r')
\braket{\vect r'}{\psi^+}\\
&=-\frac{(2\pi)^3}{4\pi}\frac{2m}{\hbar^2}
\mel{\vect k'}{\op V}{\psi^+}
\equiv f(\vect k,\vect k').
\end{aligned}\end{equation}
Here we have taken
\begin{equation}
\braket{\vect r}{\varphi}=\frac{e^{\ii\vect k\cdot\vect r}}{(2\pi)^{3/2}}.
\end{equation}
Define the transition operator $\op T$ by
\begin{equation}
\op V\ket{\psi^+}\equiv\op T\ket\varphi.
\end{equation}
Then the scattering amplitude is
\begin{equation}
\boxed{f(\vect k,\vect k')=-\frac{(2\pi)^3}{4\pi}\frac{2m}{\hbar^2}
\mel{\vect k'}{\op T}{\vect k}.}
\end{equation}
The physical meaning of the transition matrix element $\mel{\vect k'}{\op T}{\vect k}$ is clear: it is the transition amplitude from the plane wave $\ket{\vect k}$ to the free state $\ket{\vect k'}$.

If one considers the advanced Green function $G^-$, the resulting $\psi^-(\vect r)$ is a superposition of a plane wave and an inward-propagating spherical wave, rather than a solution of the outgoing scattering problem. Advanced and causal Green functions are relevant to time evolution in dynamical systems.

\subsection*{4) Optical theorem}
The forward scattering amplitude is
\begin{equation}
f(0,\varphi)=f(\vect k,\vect k)=-\frac{(2\pi)^3}{4\pi}\frac{2m}{\hbar^2}
\mel{\vect k}{\op V}{\psi^+}.
\end{equation}
From the Lippmann--Schwinger equation (with $\ket\varphi=\ket{\vect k}$),
\begin{equation}\begin{aligned}
\ket{\psi^+}&=\ket\varphi+\frac1{E-\op H_0+\ii\epsilon}\op V\ket{\psi^+},\\
\bra{\vect k}&=\bra{\psi^+}-\bra{\psi^+}\op V
\frac1{E-\op H_0-\ii\epsilon},
\end{aligned}\end{equation}
so that
\begin{equation}\begin{aligned}
\mel{\vect k}{\op V}{\psi^+}
&=\left(\bra{\psi^+}-\bra{\psi^+}\op V\frac1{E-\op H_0-\ii\epsilon}\right)\op V\ket{\psi^+}\\
&=\mel{\psi^+}{\op V}{\psi^+}
-\bra{\psi^+}\op V\frac1{E-\op H_0-\ii\epsilon}\op V\ket{\psi^+}.
\end{aligned}\end{equation}
Using
\begin{equation}
\frac1{E-\op H_0-\ii\epsilon}
=\operatorname{Pr}\frac1{E-\op H_0}+\ii\pi\delta(E-\op H_0),
\end{equation}
and noting that $\op V$ and $\op V(\operatorname{Pr}(E-\op H_0)^{-1})\op V$ are Hermitian operators, whose expectation values in $\ket{\psi^+}$ are real, we find
\begin{equation}\begin{aligned}
\operatorname{Im}\mel{\vect k}{\op V}{\psi^+}
&=-\pi\bra{\psi^+}\op V\delta(E-\op H_0)\op V\ket{\psi^+}\\
&=-\pi\int\dd^3k'\,
\mel{\psi^+}{\op V}{\vect k'}\mel{\vect k'}{\op V}{\psi^+}
\delta(E-\hbar^2k'^2/2m)\\
&=-\pi\int\dd^3k'\,
\left|\mel{\vect k'}{\op V}{\psi^+}\right|^2
\delta\left(E-\frac{\hbar^2k'^2}{2m}\right)\\
&=-\frac{\pi m k}{\hbar^2}\int\dd\Omega'\,
\left|\mel{\vect k'}{\op V}{\psi^+}\right|^2.
\end{aligned}\end{equation}
Because
\begin{equation}
\sigma(\vect k',\vect k)=|f(\vect k',\vect k)|^2
=\frac{(2\pi)^6}{16\pi^2}\frac{4m^2}{\hbar^4}
\left|\mel{\vect k'}{\op V}{\psi^+}\right|^2,\qquad
\sigma_{\mathrm{tot}}=\int\dd\Omega\,\sigma(\theta,\varphi),
\end{equation}
we obtain
\begin{equation}\begin{aligned}
\operatorname{Im}\mel{\vect k}{\op V}{\psi^+}
&=-\frac{\hbar^2k}{16\pi^3m}\int\dd\Omega'\,\sigma(\theta,\varphi)
=-\frac{\hbar^2k}{16\pi^3m}\sigma_{\mathrm{tot}}.
\end{aligned}\end{equation}
Therefore
\begin{equation}\begin{aligned}
\operatorname{Im}f(\theta=0)
&=-\frac{(2\pi)^3}{4\pi}\frac{2m}{\hbar^2}
\operatorname{Im}\mel{\vect k}{\op V}{\psi^+}\\
&=\frac{k}{4\pi}\sigma_{\mathrm{tot}}.
\end{aligned}\end{equation}
This relation between the forward scattering amplitude and the total scattering cross section is called the \emph{optical theorem}:
\begin{equation}
\boxed{\operatorname{Im}f(0)=\frac{k}{4\pi}\sigma_{\mathrm{tot}}}.
\end{equation}